\section{Experiments}
\label{sec:experiments}

% NOTE: this file is the Experiments/Setup section only. The Results section
% previously duplicated here has moved to results_section_draft.tex; keeping both
% produced duplicate table and section label definitions.

\subsection{Setup}
We evaluate MAYA in an online class-incremental setting: classes arrive sequentially in disjoint tasks and task identity is unavailable at inference. All backbones are frozen throughout; features are extracted once and cached, so every method in our comparison sees identical inputs.

\textbf{Backbones.} ResNet50 ($D{=}2048$), DINOv3 ViT-7B/16 (\texttt{vit\_7b\_patch16\_dinov3.lvd1689m}, $D{=}4096$), and SigLIP2 ViT-g-opt/16-384 ($D{=}1536$). These span a supervised convolutional encoder, a large self-supervised transformer, and a contrastively pre-trained vision-language encoder.

\textbf{Benchmarks.} ImageNet-R (200 classes of artistic and rendered styles, 20 tasks of 10 classes), TinyImageNet (200 classes, 20 tasks of 10 classes), and ObjectNet (300 classes of out-of-context object views, 15 tasks of 20 classes). Within each task we use an $80/20$ train/test split.

\textbf{Baselines.} We compare against two families, and we keep the distinction explicit because it determines what a memory comparison means.

The first family \emph{stores sample-derived data}, as MAYA does: Nearest Class Mean (NCM) on raw backbone features; Basic Replay, a feature buffer classified by nearest neighbour; ER+MLP, which trains a multi-layer perceptron head on buffered features; a memory-matched raw-vector exemplar buffer given the same budget MAYA consumes; and a Nodes-Only ablation using the episodic memory without the global projection.

The second family \emph{stores only sufficient statistics}: Deep SLDA, which maintains running class means and a shared covariance; Exact ACL, which solves recursive least squares incrementally via Sherman-Morrison updates; and RanPAC, which expands features to $M{=}10{,}000$ random dimensions before an analytic classifier. These methods retain no sample-derived vectors, and therefore offer a privacy property MAYA does not.

We do not compare against prompt-based frozen-backbone methods (L2P, DualPrompt, CODA-Prompt). These optimize inserted prompt parameters by gradient descent through the backbone, which is orthogonal to the gradient-free regime studied here and impractical at the scale of a 7B encoder. We regard them as complementary rather than competing.

\textbf{Metrics.} We report Average Incremental Accuracy (AIA), the mean of the accuracy measured after each task, and average forgetting. Alongside accuracy we track the exact memory footprint of all persistent state --- replay buffers, covariance and precision matrices, projection matrices, and episodic nodes --- measured by summing tensor sizes rather than estimated analytically.

\subsection{Implementation Details}
\label{sec:repro}

Table~\ref{tab:hparams} lists every hyperparameter. We did not tune per dataset except where noted; Section~\ref{sec:ablations} shows the analytic projection is insensitive to $\lambda$ across more than an order of magnitude.

\begin{table}[htbp]
\centering
\caption{Hyperparameters and protocol. Identical across all nine dataset $\times$ backbone configurations unless stated.}
\label{tab:hparams}
\begin{tabular}{lll}
\toprule
\textbf{Symbol} & \textbf{Meaning} & \textbf{Value} \\
\midrule
$K$        & max episodic nodes per class      & $128$ (swept over $\{1,16,64,128\}$) \\
$k_{\max}$ & $k$-means clusters per class per block & $8$ \\
$p$        & ETF / alignment dimension          & $256$ \\
$\lambda$  & RLS regularization                 & $0.1$ (swept over $\{0.01,0.1,0.5\}$) \\
$\beta$    & global/episodic blend              & selected per task on validation, grid $\{0.0,\dots,1.0\}$ \\
$T$        & node scoring temperature           & $0.08$ ($0.12$ for SigLIP2 / ResNet50) \\
--         & consolidation interval             & every block \\
--         & batch size                         & $2048$ \\
--         & train/test split per task          & $80/20$ \\
--         & seed                               & $42$ (single seed; see Section~\ref{sec:limitations}) \\
\bottomrule
\end{tabular}
\end{table}

$\beta$ is not a fixed constant. At each task it is selected on a held-out validation split drawn from the current task's training data, which is why Table~\ref{tab:components} reports both the fixed-$\beta$ endpoints and the tuned result. No test data is used for selection.

Two evaluation protocols appear in this paper and we state the difference rather than eliding it. The main comparisons (Tables~\ref{tab:mechanics},~\ref{tab:matched},~\ref{tab:sota}) shuffle each task's data before the $80/20$ split; the component ablations (Table~\ref{tab:components}) use a fixed split. Numbers are therefore comparable within a table but not across the two groups; the residual difference on ObjectNet/DINOv3 is $0.6$ AIA points.

All experiments run on cached features, so the reported cost excludes backbone forward passes, which are shared by every method and paid once per dataset.
